% ── §2  Background ────────────────────────────────────────────────────────────
\section{Background}
\label{sec:background}

\subsection{Tail Latency at Scale}
Dean and Barroso~\cite{dean2013tail} established that at scale, even a
small fraction of slow responses dominates the perceived latency of a
system: a service with 1\% of requests at 1\,s latency sees \pnn{99}
at that level regardless of the median.
This ``long tail'' effect is amplified in multi-stage pipelines because
each stage adds delay, and the slow path of each stage compounds.

\subsection{Apache Flink}
Apache Flink~\cite{carbone2015flink} is a stateful dataflow engine that
executes directed acyclic operator graphs over unbounded streams.
Key knobs relevant to tail latency:
\begin{itemize}
  \item \textbf{Checkpoint interval} ($c$): Flink snapshots operator
    state to durable storage every $c$\ms via a Chandy-Lamport barrier
    protocol. Short intervals increase fault-tolerance but add barrier
    overhead and stall operator progress.
  \item \textbf{Buffer timeout} ($b$): Flink batches output records in
    network buffers for up to $b$\ms before flushing.
    Low $b$ reduces latency; high $b$ increases throughput.
  \item \textbf{Parallelism} ($p$): the number of parallel subtasks
    per operator. Under-parallelism causes queue buildup; over-parallelism
    wastes slots and increases coordination overhead.
\end{itemize}

\subsection{Apache Kafka}
Kafka~\cite{kreps2011kafka} is the standard ingest layer for Flink
pipelines. Producer-side \texttt{linger.ms} delays sending to allow
batch accumulation; high values increase throughput but add latency to
events at the leading edge of a burst.
Consumer lag (records not yet consumed by Flink) is a leading indicator
of backpressure propagation.

\subsection{Explainable AI (XAI)}
SHAP (SHapley Additive exPlanations)~\cite{lundberg2017shap} assigns to
each feature a contribution value that satisfies the Shapley axioms
(efficiency, symmetry, dummy, additivity).
For tree-based models, \texttt{TreeExplainer} computes exact SHAP values
in $O(TLD^2)$ time, where $T$ is the number of trees, $L$ the number of
leaves, and $D$ the maximum depth---making it practical in a 10-second
control loop.

\subsection{Contextual Bandits}
A contextual bandit selects one of $K$ actions given a context vector,
observes a scalar reward, and updates a policy~\cite{sutton2018rl}.
Unlike full RL, there is no credit-assignment across time steps, which
simplifies online learning and analysis.
We use XGBoost~\cite{chen2016xgboost} as the reward predictor because:
(a)~it is natively interpretable via SHAP;
(b)~it handles tabular, non-stationary data well;
(c)~online retraining is fast ($<$\,100\ms for $<$\,10k samples).
